\documentclass[UTF8,12pt]{ctexart}
\usepackage[a4paper,margin=24mm]{geometry}
\usepackage{amsmath,amssymb,bm,graphicx,adjustbox}
\newcommand{\fitmath}[1]{\begin{adjustbox}{max width=\linewidth}$\displaystyle #1$\end{adjustbox}}
\setlength{\parindent}{0pt}
\setlength{\parskip}{.7em}
\sloppy
\allowdisplaybreaks
\begin{document}
\section*{1992 年考研数学三 · 参考答案}
\subsection*{1992年真题参考答案}

\subsection*{（试卷IV）}

\subsection*{一、填空题}

（1）\(\displaystyle (10,\allowbreak 20]\)。 （2）\(\displaystyle (0,\allowbreak 4)\)。


（3）


\[\fitmath{\displaystyle \int_0^1dx\displaystyle \int_0^{x^2}f(x,y)\,dy+\displaystyle \int_1^{\sqrt2}dx\displaystyle \int_0^{\sqrt{2-x^2}}f(x,y)\,dy.}\]


（4）\(\displaystyle (-1)^{mn}ab\)。 （5）\(\displaystyle \dfrac{\displaystyle 1}{\displaystyle 1260}\)。


\subsection*{二、选择题}

（1）B。 （2）D。 （3）A。 （4）B。 （5）C。


三、函数 \(\displaystyle f(x)\) 在 \(\displaystyle x=1\) 处不连续。若修改定义，令 \(\displaystyle f(1)=-\dfrac{\displaystyle 4}{\displaystyle \pi^2}\)，则函数 \(\displaystyle f(x)\) 在 \(\displaystyle x=1\) 处连续。


四、\(\displaystyle -e^{-x}\operatorname{arccot}e^x-x+\dfrac{\displaystyle 1}{\displaystyle 2}\ln(1+e^{2x})+C\)，其中 \(\displaystyle C\) 为任意常数。


五、\(\displaystyle \dfrac{\displaystyle \partial^2z}{\displaystyle \partial x\partial y}=\cos(xy)-xy\sin(xy)-\dfrac{\displaystyle 1}{\displaystyle y^2}\varphi'_v-\dfrac{\displaystyle x}{\displaystyle y^2}\varphi''_{uv}-\dfrac{\displaystyle x}{\displaystyle y^3}\varphi''_{vv}\)。


六、\(\displaystyle f(x)=\dfrac{\displaystyle 1}{\displaystyle 2}e^{-2x}+x-\dfrac{\displaystyle 1}{\displaystyle 2}\)。


七、证明略（可考虑函数 \(\displaystyle f(x)=\arctan x-\dfrac{\displaystyle 1}{\displaystyle 2}\arccos\dfrac{\displaystyle 2x}{\displaystyle 1+x^2}-\dfrac{\displaystyle \pi}{\displaystyle 4}\)，证明 \(\displaystyle f\prime(x)=0\)）。


八、（1）\(\displaystyle V(\xi)=\dfrac{\displaystyle \pi}{\displaystyle 2}(1-e^{-2\xi})\)，\(\displaystyle a=\dfrac{\displaystyle 1}{\displaystyle 2}\ln2\)。（2）所求切点为 \(\displaystyle (1,\allowbreak 1/e)\)，最大面积为 \(\displaystyle 2/e\)。


九、（1）\(\displaystyle x=0,\allowbreak y=-2\)。（2）


\[\fitmath{\displaystyle P=\begin{pmatrix}0&0&1\\2&1&0\\-1&1&-1\end{pmatrix}.}\]


十、（1）\(\displaystyle \lambda=1\)。（2）证明略（反证法）。


十一、\(\displaystyle C\) 是正定矩阵。


十二、\(\displaystyle \alpha=1-0.95^{100}-100\times0.05\times0.95^{99}-\dfrac{\displaystyle 100\times99}{\displaystyle 2}\times0.05^2\times0.95^{98}\)，利用泊松分布，\(\displaystyle \alpha\approx0.87\)。


十三、\(\displaystyle P\{X=0\}=0.504,\allowbreak \ P\{X=1\}=0.398,\allowbreak \ P\{X=2\}=0.092,\allowbreak \ P\{X=3\}=0.006\)，\(\displaystyle E(X)=0.6,\allowbreak D(X)=0.46\)。


十四、（1）\(\displaystyle f_X(x)=\begin{cases}e^{-x},\allowbreak &x>0,\allowbreak \\0,\allowbreak &x\le0.\end{cases}\) （2）\(\displaystyle P\{X+Y\le1\}=1+e^{-1}-2e^{-1/2}\)。


\subsection*{（试卷V）}

\subsection*{一、填空题}

（1）\(\displaystyle (2t+1)e^{2t}\)。 （2）【同试卷IV 第一、（1）题】 （3）\(\displaystyle \arcsin(1-x^2)\)；\(\displaystyle [-\sqrt2,\allowbreak \sqrt2]\)。


一、（4）\(\displaystyle 4\)。 （5）\(\displaystyle \dfrac{\displaystyle 5}{\displaystyle 8}\)。


\subsection*{二、选择题}

（1）【同试卷IV 第二、（1）题】 （2）D。 （3）C。 （4）B。 （5）【同试卷IV 第二、（4）题】


三、\(\displaystyle -\dfrac{\displaystyle 4}{\displaystyle \pi^2}\)。


四、【同试卷IV 第四题】


五、\(\displaystyle f(x)=\cos x-x\sin x+C\)，其中 \(\displaystyle C\) 为任意常数。


六、【同试卷IV 第五题】


七、（1）总利润函数 \(\displaystyle R-C=-10+72x+15x^2-x^3\)。（2）当产量为12时，总利润最大。


八、证明略（可考虑函数 \(\displaystyle f(x)=x+p+q\cos x\)，利用介值定理）。


九、（1）所求切线方程为 \(\displaystyle y-\dfrac{\displaystyle 1}{\displaystyle x_0^2}=-\dfrac{\displaystyle 2}{\displaystyle x_0^3}(x-x_0)\)。（2）所求长度为 \(\displaystyle \dfrac{\displaystyle 3\sqrt3}{\displaystyle 2}\)。


十、


\[\fitmath{\displaystyle X=\begin{pmatrix}2&0&1\\0&3&0\\1&0&2\end{pmatrix}.}\]


十一、\(\displaystyle \lambda=1\)。 十二、\(\displaystyle |A|=1\)。


十三、【同试卷IV 第十二题】 十四、【同试卷IV 第三题】

\end{document}
